\documentclass[justified, nobib]{tufte-handout}

\title{Notes: hyperbolic conservation laws \& method of characteristics}

\author[Wojciech Sadowski]{Wojciech Sadowski}

%\date{28 March 2010} % without \date command, current date is supplied

% \geometry{showframe} % display margins for debugging page layout

\usepackage{fontspec}
% \setmainfont[Renderer=Basic, Numbers=OldStyle, Scale = 1.0]{TeX Gyre Pagella}
% \setsansfont[Renderer=Basic, Scale=0.90]{TeX Gyre Heros}
% \setmonofont[Renderer=Basic]{TeX Gyre Cursor}

% Set up the spacing using fontspec features
\renewcommand\allcapsspacing[1]{{\addfontfeature{LetterSpace=15}#1}}
\renewcommand\smallcapsspacing[1]{{\addfontfeature{LetterSpace=10}#1}}

% setting ptions for included pictures
\usepackage{graphicx} % allow embedded images
\graphicspath{{graphics/}} % set of paths to search for images
\setkeys{Gin}{
  width=\linewidth,
  totalheight=\textheight,
  keepaspectratio
}

\hypersetup{colorlinks}

\usepackage[compatibility=false, font=footnotesize]{caption}
% * options are necessary to fix the bug that autoref was causing,
%   figure was constantly changed to section
\DeclareCaptionFont{blue}{\color{blue}}
\captionsetup{labelfont={blue}}

% color and color palettes
\usepackage{xcolor}

% drawings collor palette
\definecolor{p11}{HTML}{33cfff}
\definecolor{p12}{HTML}{ffa733}
\definecolor{p13}{HTML}{ff4133}


% drawing inside tex
\usepackage{tikz} 
\usetikzlibrary{decorations.markings}
\usetikzlibrary{calc}
\usetikzlibrary{shapes}

\usepackage{pgfplots}
\pgfplotsset{compat=newest}
\usepgfplotslibrary{fillbetween}
\usepackage{xcolor}

% math things
\usepackage{
  amsmath,  % extended mathematics
  amsthm,
  siunitx,  % non-stacked fractions and better unit spacing
  bm,       % bold math symbols
  physics2,  % a set of usefull symbols for differentiation 
  fixdif,
  derivative
}

\usephysicsmodule{ab, nabla.legacy}

% table things
\usepackage{
  booktabs, % book-quality tables
  multicol  % multiple column layout facilities
} 


% writing references
\usepackage{cleveref}  % must be loded after amsmath

% quoting text and friends
\usepackage{fancyvrb}         % extended verbatim environments
\fvset{fontsize=\normalsize}  % default font size for fancy-verbatim environments
\usepackage{
  dirtytalk, % provides \say command for easy quiting 
  nth       % provides \nth command for 1-st, 2-nd, etc
}        


% citing sources
\usepackage{natbib}
\setcitestyle{authoryear}
\bibliographystyle{plainnat}

% misc 
\usepackage{lipsum}   % filler text

\input{../src/commands.tex}
\newtheorem{theorem}{Theorem}
\newtheorem{lemma}{Lemma}

\newcolumntype{L}[2]{>{\raggedleft#2}p{#1\textwidth}}

\input{../src/symbols.tex}

\newcommand{\Q}[1]{\subsection{Q:~{#1}}}


\begin{document}

\maketitle% this prints the handout title, author, and date

\tableofcontents
% \begin{abstract}
% \noindent
%   \lipsum[1]{}
% \end{abstract}
\section{One-dimensional advection equation}
Given a one-dimensional advection equation 
\begin{equation}
  \pdv{u}{t} + \pdv{au}{x} = 0, \quad a = \const{}
  \label{eq:linear-pde}
\end{equation}

\Q{Define a type of this PDE}
General form of \nth{2} order PDE is given as 
\[
  A\pdv[order=2]{u}{x} + \pdv{u}{x,y}
  + C\pdv[order=2]{u}{y} 
  + D\pdv{u}{x} + E\pdv{u}{y} + Fu + G = 0.
\]
The equation is
\begin{description}
  \item[\ul{hyperbolic}] if \(B^2 - AC > 0\)
  \item[\ul{parabolic}] if \(B^2 - AC = 0\)
  \item[\ul{elliptic}] if \(B^2 - AC < 0\)
\end{description}
This is one of the methods of determining that by the way, but not \emph{the
only one.}

In our case, the equation is a \nth{1} order PDE, so we have to increase its order
by differentiating~\eqref{eq:linear-pde} with respect to time and space:
\[
  \pdv[order=2]{u}{t} + \pdv{au}{x,t} = 0,
\]
\[
  \pdv{u}{x,t} + \pdv[order=2]{au}{x} = 0.
\]
We can isolate the term \(\pdv{u}/{x,t}\) from the second equation and substitute 
it into the first, to get the following \nth{2} order equation
\[
  \pdv[order=2]{u}{t} - \pdv[order=2]{au}{x} = 0,
\]
which is hyperbolic (\(B = 0\), \(A=1\) and \(C=-a\)).\footnote{
  There is another way of doing that, described in the \say{review} slides. Try
  it out. Send me the scans if you want to share it with others as a part of
  this note.}

\Q{How is this form of the equation called?}
This is the conservation law, describing the evolution of a variable \(u\). It is
written in a conservative (strong) form. We can rewrite it introducing a flux 
function \(f=f(u)\), so that 
\[
  \pdv{u}{t} + \pdv{f(u)}{x} = 0, \quad f(u) = au
\]

\Q{Quasi-linear form}
We are trying to achieve an equation of the following form
\[
  \pdv{u}{t} + \mathrm{something}\times\pdv{u}{x} = 0,
\]
where \say{something} is \emph{evaluated} independently of \(u\) (in practice,
in numerical methods, it is evaluated according to the best approximation of
\(u\) available in a given moment). Then PDE will be linear. 

Linearization of the flux function involves computing a derivative
\[
  \pdv{f\ab(u(x))}{x} = \pdv{f(u)}{u}\pdv{u}{x},
\]
so at the end we have
\[
  \pdv{u}{t} + \pdv{f(u)}{u}\pdv{u}{x} = 0.
\]
Assuming that \(f(u) = au\), then \(\pdv{f}{u} = a\) and the equation becomes
\[
  \pdv{u}{t} + a\pdv{u}{x} = 0.
\]
This is of course a trivial example as \(a=\const\), so the computation of the 
derivative is simply not necessary.



\section{Method of characteristics}


% \begin{fullwidth}
% \bibliography{refs}
% \end{fullwidth}


\end{document}
